\section{Analytic-arc valuation}\label{sec:arc1}

Let \(v(\eta)=(a(\eta),b(\eta),c(\eta),d(\eta))\) be an analytic germ and set
\[
\alpha=\ord_\eta\ell(v(\eta)),\qquad
\beta=\ord_\eta Q(v(\eta)),
\]
with order \(\infty\) for the zero germ.  The three-scale pullback is
\(\e=c_\e s^b\), \(t=c_t s^a\), \(\eta=c_\eta s^c\), with
positive exponents \(a,b,c\) and non-zero constants.

\begin{theorem}[Analytic-arc lower hull]\label{thm:arc}
The only candidate lower points are
\[
P_\ell=(1,3,\alpha),\qquad P_Q=(2,4,\beta),
\]
omitting a point whose order is \(\infty\).  Their pulled costs are
\[
\kappa_\ell=b+3a+c\alpha,\qquad
\kappa_Q=2b+4a+c\beta.
\]
The lower hull is classified as follows:
\begin{enumerate}[label=(\roman*)]
\item if \(\alpha\le\beta<\infty\), \(P_\ell\) dominates;
\item if \(\beta<\alpha<\infty\), both points are vertices, with a tie on
      \(c(\alpha-\beta)=a+b\);
\item if \(\alpha=\infty\) and \(\beta<\infty\), only \(P_Q\) remains;
\item if \(\alpha=\beta=\infty\), the arc is exactly dark.
\end{enumerate}
\end{theorem}
\begin{proof}
The endpoint theorem gives \(q\ge p+2\) for every positive perturbation degree
\(p\); the \(p=0\) layer is zero by baseline darkness.  Every positive-degree
coefficient polynomial vanishes on the exact-dark locus.  Since
\((L,Q)\) is radical by \cref{prop:darkideal}, each such coefficient belongs to
\((L,Q)\).  More sharply, complete first-order protection shows that every
\(p=1\) coefficient lies in \((L)\).  After restriction to the analytic arc,
a \(p=1\) term therefore has order at least \(\alpha\) and cost at least
\(b+3a+c\alpha=\kappa_\ell\).

For \(p\ge2\), coefficient membership in \((L,Q)\) gives order at least
\(\min(\alpha,\beta)\), while endpoint support gives cost at least
\[
pb+(p+2)a+c\min(\alpha,\beta).
\]
If \(\alpha\le\beta\), this exceeds \(\kappa_\ell\) by at least \(a+b\).
If \(\beta<\alpha\), it is at least \(\kappa_Q\), with equality possible only
at the displayed \((2,4,\beta)\) boundary candidate.  Thus no hidden
higher-order coefficient can lie below the two stated candidates.  Comparing
their costs gives
\(\kappa_\ell-\kappa_Q=c(\alpha-\beta)-(a+b)\), which yields the four cases.
If both germs vanish, every positive-degree coefficient vanishes by membership
in \((L,Q)\), while the baseline layer is already dark; hence the whole arc is
exactly dark.
\end{proof}

On a tie, write \(\ell=\ell_\alpha\eta^\alpha+\cdots\) and
\(Q=Q_\beta\eta^\beta+\cdots\).  The candidate amplitude coefficient is
\begin{equation}\label{eq:tiecoef}
\frac{i\ell_\alpha c_\e c_\eta^\alpha c_t^3}{6}
+\frac{Q_\beta c_\e^2c_\eta^\beta c_t^4}{12}.
\end{equation}

\begin{proposition}[Field-corrected tie rule]\label{prop:tiefield}
Over \(\C\), the tied coefficient can cancel at
\begin{equation}\label{eq:ctcancel}
c_t=-\frac{2i\ell_\alpha c_\eta^{\alpha-\beta}}{Q_\beta c_\e}.
\end{equation}
For a real-analytic Hermitian family with real
\(\ell_\alpha,Q_\beta,c_\e,c_\eta,c_t\ne0\), the two terms in
\cref{eq:tiecoef} have orthogonal real and imaginary parts and cannot cancel.
No real non-cancellation statement is made for complex leading data.
\end{proposition}
\begin{proof}
Factor the non-zero cubic term from \cref{eq:tiecoef}; the remaining affine
factor in \(c_t\) gives \cref{eq:ctcancel}.  Under the stated real hypotheses,
the first summand is purely imaginary and the second purely real.
\end{proof}

\begin{figure}[ht]
\centering\begin{tikzpicture}[scale=.9,>=Latex]
\draw[->] (0,0)--(5,0) node[right] {$p$};
\draw[->] (0,0)--(0,4.4) node[above] {$q$};
\draw[->] (0,0)--(-1.8,-1.4) node[below left] {$\eta$-order};
\coordinate (L) at (1.1,3.1); \coordinate (Q) at (3.4,2.0);
\draw[very thick,blue!60!black] (L)--(Q);
\fill[blue!70!black] (L) circle (2.6pt); \fill[orange!80!black] (Q) circle (2.6pt);
\node[above left,font=\small] at (L) {$P_\ell=(1,3,\alpha)$};
\node[below right,font=\small] at (Q) {$P_Q=(2,4,\beta)$};
\draw[red!70!black,thick,dashed] (-.3,3.85)--(4.35,1.15);
\node[red!70!black,font=\small,rotate=-30] at (3.3,3.15) {tie weight plane};
\end{tikzpicture}

\caption{The two analytic-arc candidates.  A weight plane can expose one point
or their joining edge; coefficient cancellation is a separate test on that
edge.}
\label{fig:arclower}
\end{figure}

\subsection{A tie and its next survivor}

Take \(v(\eta)=(1+\eta^2,0,-1,0)\), so \(\alpha=2\), \(\beta=0\), and use
equal scale exponents.  At total order six,
\[
[s^6]K=\frac{c_\e^2c_t^4}{12}+\frac{i c_\e c_t^3}{6}.
\]
With \(c_\e=1,c_t=-2i\), this cancels.  Exact expansion gives
\[
[s^7]K=\frac{c_\e c_t^4}{12}-\frac{i c_\e^2c_t^5}{60}=\frac45.
\]
The survivor \(4/5\) is \evidence{VERIFIED-EXACT}.  Its value cannot be inferred
from \((\alpha,\beta)\) alone; it uses later coefficients of the pulled germ.

Regular analytic reparametrisation preserves \(\alpha,\beta\) and leading
scales.  A singular reparametrisation \(\eta=\xi^r\) multiplies both finite
orders by \(r\).  Analyticity is essential: the smooth-flat germ
\(e^{-1/\eta^2}\) has zero Taylor series at the origin without vanishing on a
punctured neighbourhood.

Finally, if an analytic arc is dark, the product factorisation on \(L=0\) and
the integral-domain property of convergent power-series rings show that the
whole germ lies in one of the two dark components.  It cannot switch components
infinitely often near the origin while remaining analytic.
